\subsection{Primary decomposition and extension of scalars}

In this no., A and B will denote two rings and we shall consider a ring homomorphism \(\rho: \mathbf{A} \to \mathbf{B}\) which makes B into an A-algebra; recall that, for every B-module F, \(\rho_*(\mathbf{F})\) is the commutative group F with the A-module structure defined by \(a.y = \rho(a)y\) for all \(a \in \mathbf{A}, y \in \mathbf{F}\) .

LEMMA 1. Let A be a Noetherian ring, p a prime ideal of A, E an A-module whose annihilator contains a power of p and such that \(\operatorname{Ass}(\mathbf{E}) = \{\mathfrak{p}\}\) and F a B-module such that \(\rho_{*}(\mathbf{F})\) is a flat A-module. The condition \(\mathfrak{P} \in \operatorname{Ass}_{\mathbf{B}}(\mathbf{E} \otimes_{\mathbf{A}} \mathbf{F})\) then implies \(\overline{\rho}^{1}(\mathfrak{P}) = \mathfrak{p}\) .

If \(n\) is such that \(p^n\mathbf{E} = 0\) , then \(p^n\mathbf{B} \subset \operatorname{Ann}(\mathbf{E} \otimes_{\mathbf{A}} \mathbf{F})\) , whence \(p^n\mathbf{B} \subset \mathfrak{P}\) , which implies \(p^n \subset \overline{\rho}^1(\mathfrak{P})\) and therefore \(p \subset \overline{\rho}^1(\mathfrak{P})\) since \(\overline{\rho}^1(\mathfrak{P})\) is prime. Moreover, if \(a \in \mathbf{A} - p\) , the homothety \(h\) with ratio \(a\) on \(\mathbf{E}\) is injective (§ 1, no. 1, Corollary 2 to Proposition 2); as \(h \otimes l_{\mathbf{F}}\) is the homothety \(h'\) with ratio \(\rho(a)\) on \(\mathbf{E} \otimes_{\mathbf{A}} \mathbf{F}\) and \(\rho_*(\mathbf{F})\) is flat, \(h'\) is injective (Chapter I, § 2, no. 2, Definition 1); this proves that \(\rho(a) \notin \mathfrak{P}\) , whence \(\overline{\rho}^1(\mathfrak{P}) = p\) .

THEOREM 2. Let \(\rho: \mathbf{A} \to \mathbf{B}\) be a ring homomorphism, E an A-module and F a B-module such that \(\rho_*(\mathbf{F})\) is a flat A-module. Then

\[
\operatorname{Ass} _ {\mathbf {B}} (\mathrm{E} \otimes_ {\mathbf {A}} \mathrm{F}) \supset \bigcup_ {\mathfrak {p} \in \operatorname{Ass} _ {\mathbf {A}} (\mathrm{E})} \operatorname{Ass} _ {\mathbf {B}} (\mathrm{F} / \mathfrak {p} \mathrm{F}).\tag{5}
\]

When A is Noetherian, the two sides of (5) are equal.

Let \(\mathfrak{p} \in \operatorname{Ass}_{\mathbf{A}}(\mathbf{E})\) ; by definition there exists an exact sequence

\[
0 \rightarrow \mathrm{A} / \mathfrak {p} \rightarrow \mathrm{E}.
\]

Since F is a flat A-module we derive an exact sequence

\[
0 \rightarrow \mathrm{F} / \mathfrak {p} \mathrm{F} \rightarrow \mathrm{E} \otimes_ {\mathrm{A}} \mathrm{F}
\]

whence \(\operatorname{Ass}_{\mathbf{B}}(\mathbf{F} / p\mathbf{F})\subset \operatorname{Ass}_{\mathbf{B}}(\mathbf{E}\otimes_{\mathbf{A}}\mathbf{F})\) , which proves the inclusion (5).

Suppose now that A is Noetherian and let us prove the opposite inclusion.

We proceed in stages:

\begin{enumerate}
\def\labelenumi{(\roman{enumi})}
\tightlist
\item
  Suppose first that E is a finitely generated A-module and that \(\operatorname{Ass}_{\mathbf{A}}(\mathbf{E})\) is reduced to a single element \(p\) . By § 1, no. 4, Theorem 1 there exists a composition series \((\mathrm{E}_i)_{0 \leqslant i \leqslant n}\) of E such that \(\mathrm{E}_i / \mathrm{E}_{i+1}\) is isomorphic to \(\mathbf{A} / p_i\) , where \(p_i\) is a prime ideal of A; moreover (§ 1, no. 4, Theorem 2 and no. 3, Proposition 7) all the \(p_i\) contain \(p\) . As F is a flat A-module, the \(\mathrm{E}_i \otimes_{\mathbf{A}} \mathbf{F}\) form a composition series of \(\mathbf{E} \otimes_{\mathbf{A}} \mathbf{F}\) and \((\mathrm{E}_i \otimes_{\mathbf{A}} \mathbf{F}) / (\mathrm{E}_{i+1} \otimes_{\mathbf{A}} \mathbf{F})\) is identified with
\end{enumerate}

\[
(\mathbf {A} / \mathfrak {p} _ {i}) \otimes_ {\mathbf {A}} \mathbf {F} = \mathbf {F} / \mathfrak {p} _ {i} \mathbf {F}.
\]

Then by virtue of § 1, no. 1, Proposition 3

\[
\operatorname{Ass} _ {\mathbf {B}} (\mathrm{E} \otimes_ {\mathbf {A}} \mathrm{F}) \subset \bigcup_ {i = 0} ^ {n - 1} \operatorname{Ass} _ {\mathbf {B}} (\mathrm{F} / p _ {i} \mathrm{F}).
\]

We know that E is annihilated by a power of p (no. 1, Remark); Lemma 1 then shows that, for all \(\mathfrak{P} \in \operatorname{Ass}_{\mathbf{B}}(\mathbf{E} \otimes_{\mathbf{A}} \mathbf{F})\) , \(\overline{\rho}^{1}(\mathfrak{P}) = \mathfrak{p}\) . As \(\mathbf{F} / \mathfrak{p}_i \mathbf{F}\) is isomorphic to \((\mathbf{A} / \mathfrak{p}_i) \otimes_{\mathbf{A}} \mathbf{F}\) , \(\overline{\rho}^{1}(\mathfrak{P}') = \mathfrak{p}_i\) for all \(\mathfrak{P}' \in \operatorname{Ass}_{\mathbf{B}}(\mathbf{F} / \mathfrak{p}_i \mathbf{F})\) by Lemma 1, whence \(\operatorname{Ass}_{\mathbf{B}}(\mathbf{E} \otimes_{\mathbf{A}} \mathbf{F}) \cap \operatorname{Ass}(\mathbf{F} / \mathfrak{p}_i \mathbf{F}) = \varnothing\) if \(\mathfrak{p}_i \neq \mathfrak{p}\) , which proves the theorem in the case considered.

\begin{enumerate}
\def\labelenumi{(\roman{enumi})}
\setcounter{enumi}{1}
\tightlist
\item
  Suppose only that E is a finitely generated A-module. Let \(p_{i}\) ( \(1 \leqslant i \leqslant n\) ) be the elements of \(\operatorname{Ass}_{\mathbf{A}}(\mathbf{E})\) and let \(\{0\} = \bigcap_{i=1}^{n} Q_{i}\) be a corresponding reduced primary decomposition (no. 3); E is then isomorphic to a submodule of the direct sum of the \(E_{i} = E/Q_{i}\) and, as F is a flat A-module, \(E \otimes_{A} F\) is isomorphic to a submodule of the direct sum of B-modules \(E_{i} \otimes_{A} F\) . We deduce (§ 1, no. 1, Proposition 3 and Corollary 1 to Proposition 3) that
\end{enumerate}

\[
\operatorname{Ass} _ {\mathbf {B}} (\mathrm{E} \otimes_ {\mathrm{A}} \mathrm{F}) \subset \bigcup_ {i = 1} ^ {n} \operatorname{Ass} _ {\mathbf {B}} (\mathrm{E} _ {i} \otimes_ {\mathrm{A}} \mathrm{F}).
\]

But \(E_{i}\) is a finitely generated A-module such that \(\operatorname{Ass}_{\mathbf{A}}(\mathbf{E}_{i})\) is reduced to a single element \(p_{i}\) (no. 1, Definition 1). By (i), \(\operatorname{Ass}_{\mathbf{B}}(\mathbf{E}_{i} \otimes_{\mathbf{A}} \mathbf{F}) = \operatorname{Ass}_{\mathbf{B}}(\mathbf{F}/p_{i}\mathbf{F})\) , whence the theorem in this case.

\begin{enumerate}
\def\labelenumi{(\roman{enumi})}
\setcounter{enumi}{2}
\tightlist
\item
  General case. The B-module \(E \otimes_{A} F\) is the union of the submodules \(E' \otimes_{A} F\) , where \(E'\) runs through the set of finitely generated submodules of the A-module E. If P belongs to \(\operatorname{Ass}_{\mathbf{B}}(\mathbf{E} \otimes_{\mathbf{A}} \mathbf{F})\) , then there exists a finitely generated submodule \(E'\) of E such that \(\mathfrak{P} \in \operatorname{Ass}_{\mathbf{B}}(\mathbf{E}' \otimes_{\mathbf{A}} \mathbf{F})\) . By (ii), there exists \(\mathfrak{p} \in \operatorname{Ass}_{\mathbf{A}}(\mathbf{E}')\) such that \(\mathfrak{P} \in \operatorname{Ass}_{\mathbf{B}}(\mathbf{F}/\mathfrak{p}\mathbf{F})\) ; as \(\operatorname{Ass}_{\mathbf{A}}(\mathbf{E}') \subset \operatorname{Ass}_{\mathbf{A}}(\mathbf{E})\) , this completes the proof of Theorem 2.
\end{enumerate}

COROLLARY 1. If A is Noetherian and \(\mathfrak{P} \in \operatorname{Ass}_{\mathbf{B}}(\mathbf{E} \otimes_{\mathbf{A}} \mathbf{F})\) , then \(\overline{\rho}^{1}(\mathfrak{P}) \in \operatorname{Ass}_{\mathbf{A}}(\mathbf{E})\) and \(\overline{\rho}^{1}(\mathfrak{P})\) is the only prime ideal \(\mathfrak{p}\) of A such that \(\mathfrak{P} \in \operatorname{Ass}_{\mathbf{B}}(\mathbf{F}/\mathfrak{p}\mathbf{F})\) .

This follows from Theorem 2 and Lemma 1 applied to the case where \(\mathbf{E} = \mathbf{A} / \mathfrak{p}\) .

COROLLARY 2. Suppose that A and B are Noetherian and that B is a flat A-module. Let p be a prime ideal of A, Q ⊂ E a p-primary submodule and P a prime ideal of B. For Q ⊗A B to be a P-primary submodule of E ⊗A B, it is necessary and sufficient that pB be a P-primary ideal of B.

Let us apply Theorem 2 to the A-module E/Q and the B-module B; then \(\operatorname{Ass}_{\mathbf{A}}(\mathrm{E} / \mathrm{Q}) = \{\mathfrak{p}\}\) and \((\mathrm{E} / \mathrm{Q})\otimes_{\mathbf{A}}\mathbf{B}\) is isomorphic to \((\mathrm{E}\otimes_{\mathbf{A}}\mathrm{B}) / (\mathrm{Q}\otimes_{\mathbf{A}}\mathrm{B})\) and hence \(\operatorname{Ass}_{\mathbf{B}}((\mathrm{E}\otimes_{\mathbf{A}}\mathrm{B}) / (\mathrm{Q}\otimes_{\mathbf{A}}\mathrm{B})) = \operatorname{Ass}_{\mathbf{B}}(\mathrm{B} / \mathfrak{p}\mathbf{B})\) . To say that \(\mathbf{Q}\otimes_{\mathbf{A}}\mathbf{B}\) is \(\mathfrak{P}\) -primary in \(\mathrm{E}\otimes_{\mathbf{A}}\mathrm{B}\) therefore means that \(\operatorname{Ass}_{\mathbf{B}}(\mathrm{B} / \mathfrak{p}\mathrm{B})\) is reduced to \(\mathfrak{P}\) , whence the corollary.

Remark. Suppose that A and B are Noetherian. Let \(\mathfrak{P}\) be a prime ideal of B and let \(\mathfrak{p} = \overline{\rho}^{1}(\mathfrak{P})\) ; let us write \(S = A - \mathfrak{p}\) and let \(k(\mathfrak{p}) = S^{-1}(A / \mathfrak{p})\) be the field of fractions of \(A / \mathfrak{p}\) . Since \(\mathfrak{P}\) contains \(\mathfrak{pB}\) , \(\mathfrak{P} / \mathfrak{pB}\) is a prime ideal of \(B / \mathfrak{pB}\) . If \(\rho'\) is the composite homomorphism \(A \xrightarrow{\rho} B \to B / \mathfrak{pB}\) , we know that \(S^{-1}(B / \mathfrak{pB})\) is identified with the ring \((\rho'(S))^{-1}(B / \mathfrak{pB})\) and \(\mathfrak{P}' = S^{-1}(\mathfrak{P} / \mathfrak{pB})\) with an ideal of this ring (Chapter II, § 2, no. 2, Proposition 6); as \(\mathfrak{P} / \mathfrak{pB}\) does not meet \(\rho'(S)\) , \(\mathfrak{P}'\) is a prime ideal of \(S^{-1}(B / \mathfrak{pB})\) (Chapter II, § 2, no. 5, Proposition 11); moreover there are canonical isomorphisms between \(S^{-1}(B / \mathfrak{pB})\) , \(S^{-1}((A / \mathfrak{p}) \otimes_{A} B)\) and \((S^{-1}(A / \mathfrak{p})) \otimes_{A} B = k(\mathfrak{p}) \otimes_{A} B\) ; similarly \(S^{-1}(F / \mathfrak{pF})\) is canonically identified with \(k(\mathfrak{p}) \otimes_{A} F\) . This being so, under the hypotheses of Theorem 2, in order that \(\mathfrak{P} \in \operatorname{Ass}_{B}(E \otimes_{A} F)\) , it is necessary and sufficient that \(\mathfrak{p} \in \operatorname{Ass}_{A}(E)\) and

\[
\mathfrak {P} ^ {\prime} \in \operatorname{Ass} _ {k (\mathfrak {p}) \otimes_ {\mathbf {A}} \mathbf {B}} (k (\mathfrak {p}) \otimes_ {\mathbf {A}} \mathbf {F}).
\]

For by Theorem 2 and its Corollary 1, it amounts to verifying that the conditions

\[
" \mathfrak {P} \in \operatorname{Ass} _ {\mathbf {B}} (\mathrm{F} / \mathfrak {p} \mathrm{F})" \quad \text { and } \quad " \mathfrak {P} ^ {\prime} \in \operatorname{Ass} _ {k (\mathfrak {p}) \otimes_ {\mathbf {A}} \mathbf {B}} (k (\mathfrak {p}) \otimes_ {\mathbf {A}} \mathrm{F})"
\]

are equivalent; but, as B is Noetherian, this follows from § 1, no. 2, Corollary to Proposition 5 and the above identifications.

PROPOSITION 11. Suppose that A and B are Noetherian and that B is a flat A-module. Let E be an A-module and E' a submodule of E such that, for every ideal \(\mathfrak{p} \in \operatorname{Ass}_{\mathbf{A}}(\mathbf{E}/\mathbf{E}')\) , \(\mathfrak{p}\mathbf{B}\) is a prime ideal of B or equal to B. Let \(\mathbf{E}' = \bigcap_{\mathfrak{p} \in \operatorname{Ass}(\mathbf{E}/\mathbf{E}')} \mathbf{Q}(\mathfrak{p})\) be a reduced primary decomposition of \(\mathbf{E}'\) in E, \(\mathbf{Q}(\mathfrak{p})\) being \(\mathfrak{p}\) -primary for all \(\mathfrak{p} \in \operatorname{Ass}(\mathbf{E}/\mathbf{E}')\) .

\begin{enumerate}
\def\labelenumi{(\roman{enumi})}
\item
  If \(\mathfrak{p} \in \operatorname{Ass}(\mathrm{E} / \mathrm{E}')\) and \(\mathfrak{p}\mathbf{B} = \mathbf{B}\) , then \(\mathbf{Q}(\mathfrak{p}) \otimes_{\mathbf{A}} \mathbf{B} = \mathbf{E} \otimes_{\mathbf{A}} \mathbf{B}\) .
\item
  If \(\mathfrak{p} \in \operatorname{Ass}(\mathrm{E} / \mathrm{E}')\) and \(\mathfrak{pB}\) is prime, \(\mathsf{Q}(\mathfrak{p})\) is \(\mathfrak{pB}\) -primary in \(\mathbf{E} \otimes_{\mathbf{A}} \mathbf{B}\) .
\item
  If \(\Phi\) is the set of \(\mathfrak{p} \in \operatorname{Ass}(\mathrm{E} / \mathrm{E}')\) such that \(\mathfrak{pB}\) is prime, then
\end{enumerate}

\[
\mathbf {E} ^ {\prime} \otimes_ {\mathbf {A}} \mathbf {B} = \bigcap_ {\mathfrak {p} \in \Phi} (\mathbf {Q} (\mathfrak {p}) \otimes_ {\mathbf {A}} \mathbf {B})
\]

and this relation is a reduced primary decomposition of \(E^{\prime} \otimes_{A} B\) in \(E \otimes_{A} B\) .

If \(\mathfrak{pB} = \mathbf{B}\) , Theorem 2 applied to \(\mathbf{E} / \mathbf{Q}(\mathfrak{p})\) and \(\mathbf{B}\) shows that

\[
\operatorname{Ass} _ {\mathbf {B}} ((\mathrm{E} / \mathrm{Q} (\mathfrak {p})) \otimes_ {\mathbf {A}} \mathbf {B}) = \varnothing
\]

and, as B is Noetherian and is a flat A-module, we conclude (§ 1, no. 1, Corollary 1 to Proposition 2) that \(\mathbf{Q}(\mathfrak{p}) \otimes_{\mathbf{A}} \mathbf{B} = \mathbf{E} \otimes_{\mathbf{A}} \mathbf{B}\) . Assertion (ii) follows from Corollary 2 to Theorem 2, taking \(\mathfrak{P} = \mathfrak{p}\mathbf{B}\) . Finally the relation \(\mathbf{E}' \otimes_{\mathbf{A}} \mathbf{B} = \bigcap_{\mathfrak{p} \in \Phi} (\mathbf{Q}(\mathfrak{p}) \otimes_{\mathbf{A}} \mathbf{B})\) follows from the fact that B is a flat A-module (Chapter I, § 2, no. 6, Proposition 6); as \(\mathfrak{p} = \overline{\rho}^{1}(\mathfrak{p}\mathbf{B})\) for \(\mathfrak{p} \in \Phi\) (Lemma 1), \(\mathfrak{p}\mathbf{B} \neq \mathfrak{p}'\mathbf{B}\) for two distinct ideals \(\mathfrak{p}, \mathfrak{p}'\) of the set \(\Phi\) ; on the other hand,

\[
\operatorname{Ass} ((\mathbf {E} \otimes_ {\mathbf {A}} \mathbf {B}) / (\mathbf {E} ^ {\prime} \otimes_ {\mathbf {A}} \mathbf {B})) = \Phi
\]

by Theorem 2; we conclude from no. 3, Proposition 4 that

\[
\mathbf {E} ^ {\prime} \otimes_ {\mathbf {A}} \mathbf {B} = \bigcap_ {\mathfrak {p} \in \Phi} (\mathbf {Q} (\mathfrak {p}) \otimes_ {\mathbf {A}} \mathbf {B})
\]

is a reduced primary decomposition.

COROLLARY. Suppose that \(\mathfrak{pB}\) is prime for all \(\mathfrak{p} \in \operatorname{Ass}_{\mathbf{A}}(\mathrm{E}/\mathrm{E}')\) . Then, if \(\mathfrak{p}_1, \ldots, \mathfrak{p}_n\) are the minimal elements of \(\operatorname{Ass}_{\mathbf{A}}(\mathrm{E}/\mathrm{E}')\) , the \(\mathfrak{p}_i\mathbf{B}\) are minimal elements of

\[
\operatorname{Ass} _ {\mathbf {A}} ((\mathbf {E} \otimes_ {\mathbf {A}} \mathbf {B}) / (\mathbf {E} ^ {\prime} \otimes_ {\mathbf {A}} \mathbf {B})).
\]

It follows from Proposition 11 that in this case \(\mathfrak{p}_i\mathbf{B} \neq \mathfrak{p}_j\mathbf{B}\) for \(i \neq j\) .

\textbf{Examples}\par

\begin{enumerate}
\def\labelenumi{(\arabic{enumi})}
\item
  Let us take \(B = S^{-1}A\) , where S is a multiplicative subset of A; if A is Noetherian, the hypotheses of Proposition 11 are satisfied and we recover a part of Proposition 6 of no. 4.
\item
  Let A be a Noetherian ring, m an ideal of A and B the Hausdorff completion of A with respect to the m-adic topology; then B is a flat A-module and Theorem 2 may be applied with F = B; but in general the hypotheses of Proposition 11 are not satisfied for the prime ideals of A (Chapter III, § 2, Exercise 15 (b)).
\item
  Let A be a Noetherian ring and B the polynomial algebra \(A[X_{1},\ldots,X_{n}]\) ; B is Noetherian and is a free A-module and therefore flat. Also, if p is a prime ideal of A, B/pB is isomorphic to \((A/p)[X_{1},\ldots,X_{n}]\) , which is an integral domain, and hence pB is prime; the hypotheses of Proposition 11 are therefore satisfied for every A-module E and every submodule \(E'\) of E.
\item
  Let A be a finitely generated algebra over a field k, K an extension of k and \(B = A \otimes_{k} K\) the algebra over K obtained by extension of scalars; A and B are Noetherian and B is a free A-module and hence Theorem 2 may be applied to F = B. In certain cases (for example if k is algebraically closed) it can be shown that for every prime ideal p of A, pB is prime or equal to B; we shall return later to this example.
\end{enumerate}

